\section{Khung bộ tiêu chí, minh chứng và kết quả tự đánh giá chương trình đào tạo nghề Điện công nghiệp}

\subsection{Căn cứ pháp lý và phương pháp luận thiết lập bộ tiêu chí chương trình đào tạo}

Việc xây dựng bộ tiêu chí và hệ thống công cụ tự đánh giá cho chương trình đào tạo trình độ cao đẳng nghề Điện công nghiệp tại Trường Cao đẳng nghề Kỹ thuật -- Công nghiệp A được thiết lập trên nền tảng pháp lý trực tiếp là Chương III của Thông tư số 14/2024/TT-BLĐTBXH do Bộ Lao động -- Thương binh và Xã hội ban hành \parencite{bldtbxh2024}. Chương này quy định hệ thống chuẩn mực kiểm định đối với chương trình đào tạo trình độ trung cấp và cao đẳng, bao gồm bảy tiêu chí đánh giá toàn diện từ khâu xác định mục tiêu, thiết kế cấu trúc nội dung, điều kiện giảng viên, cơ sở vật chất cho đến công tác giám sát đầu ra và việc làm của người học.

Đồng thời, phương pháp luận thiết lập bộ tiêu chí bám sát quy định tại Quyết định số 1982/QĐ-TTg của Thủ tướng Chính phủ phê duyệt Khung trình độ quốc gia Việt Nam \parencite{qd1982}, trong đó xác định trình độ cao đẳng tương ứng với Bậc 5. Ở bậc trình độ này, người học sau khi tốt nghiệp nghề Điện công nghiệp bắt buộc phải có kiến thức thực tế và lý thuyết chuyên sâu về hệ thống cung cấp điện, trang bị điện máy công nghiệp, điều khiển tự động hóa; có kỹ năng phân tích, thiết kế, lắp đặt, vận hành, bảo dưỡng và xử lý sự cố trong các hệ thống điện công nghiệp phức tạp; đồng thời có năng lực làm việc độc lập, sáng tạo, chịu trách nhiệm cá nhân và trách nhiệm đối với nhóm làm việc.

Về phương pháp luận xây dựng chỉ báo, bộ tiêu chí tuân thủ chặt chẽ bốn nguyên tắc khoa học: tính chuẩn mực pháp lý, tính định lượng có thể đo lường được, tính tương thích với chuẩn kỹ năng nghề quốc gia và tính thực chứng của dữ liệu minh chứng. Đối với nghề Điện công nghiệp, một ngành kỹ thuật chịu tác động trực tiếp của quá trình tự động hóa và số hóa công nghiệp, các tiêu chí và tiêu chuẩn không dừng lại ở các yêu cầu cơ học truyền thống mà được cụ thể hóa thành các chỉ báo kỹ thuật chuyên sâu theo định hướng phát triển nguồn nhân lực kỹ thuật thế kỷ XXI \parencite{trankhanhduc2010} như năng lực lập trình bộ điều khiển logic khả trình (PLC), vận hành hệ thống giám sát và thu thập dữ liệu (SCADA), sử dụng phần mềm thiết kế điện chuyên dụng và tuân thủ các quy chuẩn nghiêm ngặt về an toàn, vệ sinh lao động điện.

\newpage
\subsection{Khung bộ tiêu chí và hệ thống chỉ báo đánh giá chương trình đào tạo nghề Điện công nghiệp}

Căn cứ quy định tại Điều 9, Điều 10 và Phụ lục 03 ban hành kèm theo Thông tư số 14/2024/TT-BLĐTBXH, hệ thống đánh giá chất lượng chương trình đào tạo trình độ cao đẳng nghề Điện công nghiệp được cấu trúc thành bảy tiêu chí với tổng số bốn mươi tiêu chuẩn. Thang điểm đánh giá áp dụng quy tắc tuyệt đối: mỗi tiêu chuẩn khi đáp ứng đầy đủ các yêu cầu theo quy định được tính 2.5 điểm; tiêu chuẩn không đáp ứng hoặc không có đủ minh chứng hợp lệ bị tính 0 điểm. Tổng số điểm chuẩn tối đa của toàn bộ hệ thống là 100 điểm.

\begin{table}[htbp]
\centering
\footnotesize
\setlength{\tabcolsep}{3.5pt}
\renewcommand{\arraystretch}{1.1}
\caption{\textbf{Khung cấu trúc 07 tiêu chí, 40 tiêu chuẩn kiểm định chương trình đào tạo nghề Điện công nghiệp}}
\label{tab:bang2_1}
\vspace{0.2cm}
\begin{tabularx}{\textwidth}{c l c c c X}
\toprule
\textbf{Mã TC} & \textbf{Tên tiêu chí kiểm định} & \textbf{Số TC} & \textbf{Điểm} & \textbf{Tỷ lệ} & \textbf{Yêu cầu và chỉ báo kỹ thuật trọng tâm} \\
\midrule
TC-01 & Mục tiêu, chuẩn đầu ra chương trình & 3 & 7.5 & 7.5\% & Mục tiêu phù hợp sứ mạng; chuẩn đầu ra tương thích Bậc 5 VQF; tích hợp chuẩn năng lực số ứng dụng. \\
TC-02 & Cấu trúc, nội dung chương trình, giáo trình & 7 & 17.5 & 17.5\% & Thực hành $\ge$ 65\%; doanh nghiệp tham gia rà soát; giáo trình được thẩm định, có bản quyền phần mềm. \\
TC-03 & Hoạt động đào tạo & 7 & 17.5 & 17.5\% & Tiến độ đào tạo khả thi; áp dụng dạy học dự án; đánh giá năng lực thực hành dựa trên sản phẩm thật. \\
TC-04 & Đội ngũ nhà giáo, cán bộ quản lý & 6 & 15.0 & 15.0\% & Nhà giáo đạt chuẩn chuyên môn, nghiệp vụ; $\ge$ 50\% đạt kỹ năng nghề cao; định kỳ đi thực tế doanh nghiệp. \\
TC-05 & Cơ sở vật chất, thiết bị và học liệu & 7 & 17.5 & 17.5\% & Xưởng đạt chuẩn diện tích; thiết bị đồng bộ, $\le$ 3 sinh viên/mô hình; hệ thống tiếp địa an toàn điện đạt chuẩn. \\
TC-06 & Người học và hoạt động hỗ trợ & 5 & 12.5 & 12.5\% & Tuyển sinh minh bạch; chính sách học bổng, miễn giảm học phí; câu lạc bộ kỹ thuật và hỗ trợ khởi nghiệp. \\
TC-07 & Giám sát, đánh giá chất lượng & 5 & 12.5 & 12.5\% & Duy trì sĩ số $\ge$ 85\%; việc làm sau 12 tháng $\ge$ 75\%; khảo sát tự động sự hài lòng của doanh nghiệp. \\
\midrule
\multicolumn{2}{l}{\textbf{Tổng số tiêu chuẩn kiểm định}} & \textbf{40} & \textbf{100} & \textbf{100\%} & \textbf{Đánh giá toàn diện chương trình đào tạo} \\
\bottomrule
\end{tabularx}
\end{table}

Trong cơ cấu phân bổ điểm số, nhóm tiêu chí về chuyên môn trực tiếp bao gồm Tiêu chí 2 (Cấu trúc chương trình), Tiêu chí 3 (Hoạt động đào tạo) và Tiêu chí 5 (Cơ sở vật chất) chiếm tỷ trọng lớn nhất với 17.5 điểm mỗi tiêu chí. Điều này phản ánh rõ nét triết lý của giáo dục nghề nghiệp là lấy năng lực thực hành, tính công nghệ của môi trường đào tạo và phương pháp tổ chức rèn luyện kỹ năng làm trọng tâm cốt lõi.

\subsection{Khung danh mục minh chứng chuẩn hóa cho chương trình đào tạo nghề Điện công nghiệp}

Một hệ thống tiêu chuẩn kiểm định chỉ có giá trị thực tiễn khi được bảo đảm bằng một khung danh mục minh chứng chặt chẽ, đầy đủ và có khả năng kiểm chứng độc lập. Nhằm chấm dứt tình trạng lưu trữ minh chứng phân tán, hình thức và e ngại chia sẻ số liệu từng xảy ra trong đợt tự đánh giá cấp trường năm 2026, nhà trường đã thiết lập quy tắc mã hóa minh chứng chuẩn mực theo cấu trúc mã bốn trường thông tin: ký hiệu H1 đại diện cho Báo cáo tự đánh giá chương trình đào tạo, tiếp theo là số thứ tự tiêu chí gồm hai chữ số, số thứ tự tiêu chuẩn gồm hai chữ số và số thứ tự của minh chứng thành phần gồm hai chữ số.

\begin{footnotesize}
\setlength{\tabcolsep}{3.5pt}
\renewcommand{\arraystretch}{1.15}
\begin{longtable}{>{\raggedright\arraybackslash\ttfamily}p{2.3cm} >{\raggedright\arraybackslash}p{5.5cm} >{\raggedright\arraybackslash}p{4.4cm} >{\raggedright\arraybackslash}p{2.2cm}}
\caption{\textbf{Khung danh mục minh chứng chuẩn hóa cốt lõi cho chương trình đào tạo nghề Điện công nghiệp}}\label{tab:bang2_2} \\
\toprule
\textbf{Mã MC} & \textbf{Tên danh mục minh chứng cốt lõi} & \textbf{Nguồn dữ liệu và Nơi lưu trữ} & \textbf{Đơn vị chủ trì} \\
\midrule
\endfirsthead

\multicolumn{4}{l}{\textit{Bảng~\ref{tab:bang2_2} (tiếp theo): Khung danh mục minh chứng chuẩn hóa cốt lõi}} \\
\toprule
\textbf{Mã MC} & \textbf{Tên danh mục minh chứng cốt lõi} & \textbf{Nguồn dữ liệu và Nơi lưu trữ} & \textbf{Đơn vị chủ trì} \\
\midrule
\endhead

\midrule
\multicolumn{4}{r}{\textit{(Xem tiếp trang sau)}} \\
\endfoot

\bottomrule
\endlastfoot

H1.01.01.01 & Quyết định ban hành mục tiêu chương trình đào tạo; Biên bản lấy ý kiến chuyên gia doanh nghiệp. & Hồ sơ chương trình đào tạo; Văn phòng Khoa Điện & Khoa Điện \\
H1.01.02.01 & Bản mô tả chuẩn đầu ra đối sánh Bậc 5 Khung trình độ quốc gia; Quyết định công bố công khai. & Trang thông tin trường; Phòng Đào tạo & Phòng Đào tạo \\
H1.01.03.01 & Đề cương chi tiết mô-đun tích hợp năng lực số; Ma trận chuẩn đầu ra kỹ năng số nghề Điện. & Hệ thống quản lý học tập LMS; Phòng ĐBCL & Khoa Điện \\
H1.02.02.01 & Biên bản rà soát chương trình đào tạo; Văn bản góp ý của 05 doanh nghiệp đối tác ngành Điện. & Hồ sơ đào tạo; Thỏa thuận hợp tác doanh nghiệp & Phòng Đào tạo, Khoa Điện \\
H1.02.05.01 & Quyết định ban hành giáo trình; Biên bản nghiệm thu; Giấy lưu chiểu tại Thư viện trường. & Thư viện trường; Hồ sơ nghiên cứu khoa học & Thư viện, Khoa Điện \\
H1.03.01.01 & Tiến độ đào tạo toàn khóa và thời khóa biểu chi tiết; Sổ theo dõi tiến độ giảng dạy. & Phần mềm quản lý đào tạo trực tuyến & Phòng Đào tạo, Khoa Điện \\
H1.03.04.01 & Ngân hàng câu hỏi, đề thi lý thuyết và thực hành kết thúc mô-đun có đáp án chi tiết. & Phòng Khảo thí và Đảm bảo chất lượng (niêm phong) & Phòng KT\&ĐBCL \\
H1.04.02.01 & Hồ sơ văn bằng, chứng chỉ sư phạm và chứng chỉ kỹ năng nghề quốc gia của nhà giáo. & Phòng Tổ chức Cán bộ; Hồ sơ nhân sự giảng viên & Phòng Tổ chức Cán bộ \\
H1.05.01.01 & Bảng kê diện tích xưởng trang bị điện, PLC; Bản vẽ mặt bằng và tải trọng xưởng máy. & Hồ sơ kỹ thuật mặt bằng; Phòng Quản trị Thiết bị & Phòng Quản trị TB \\
H1.05.03.01 & Nhật ký bảo dưỡng máy móc; Biên bản kiểm tra an toàn điện và tiếp địa định kỳ. & Sổ nhật ký hiện trường xưởng thực hành & Khoa Điện, Phòng QTTB \\
H1.07.02.01 & Báo cáo kết quả khảo sát việc làm sinh viên tốt nghiệp sau 12 tháng (kèm số liệu đối soát). & Cơ sở dữ liệu việc làm sinh viên tốt nghiệp & Trung tâm HTDN \\
H1.07.03.01 & Báo cáo tổng hợp ý kiến phản hồi của các doanh nghiệp sử dụng lao động về nhân sự. & Báo cáo khảo sát trực tuyến; Phòng ĐBCL & Phòng Đảm bảo chất lượng \\
\end{longtable}
\end{footnotesize}

\subsection{Kết quả thẩm định và phân tích minh chứng thử cho năm tiêu chuẩn điển hình}

Để đánh giá mức độ sẵn sàng và độ tin cậy của nguồn dữ liệu trước khi đăng ký kiểm định ngoài, Hội đồng tự đánh giá nhà trường đã tiến hành thu thập thử và thẩm định độc lập hồ sơ minh chứng của năm tiêu chuẩn đại diện cho các nhóm hoạt động chuyên môn trọng yếu của chương trình đào tạo nghề Điện công nghiệp.

\begin{table}[htbp]
\centering
\footnotesize
\setlength{\tabcolsep}{3.5pt}
\renewcommand{\arraystretch}{1.15}
\caption{\textbf{Thẩm định thử hồ sơ minh chứng cho 05 tiêu chuẩn điển hình của chương trình đào tạo}}
\label{tab:bang2_3}
\vspace{0.2cm}
\begin{tabularx}{\textwidth}{l X X l}
\toprule
\textbf{Tiêu chuẩn} & \textbf{Danh mục minh chứng thử nghiệm} & \textbf{Đánh giá 04 thuộc tính chất lượng minh chứng} & \textbf{Kết luận} \\
\midrule
\textbf{TC 1.2} & Bản mô tả chuẩn đầu ra (\texttt{H1.01.02.01}); Quyết định phê duyệt; Ảnh chụp website. & Pháp lý: Tốt; Đầy đủ: Đạt; Tin cậy: Cao; Tương thích: Hoàn toàn phù hợp Khung trình độ quốc gia Bậc 5. & \textbf{Đạt yêu cầu} \\
\textbf{TC 1.3} & Đề cương mô-đun CAD điện, PLC (\texttt{H1.01.03.01}); Phiếu chấm bài tập mô phỏng. & Pháp lý: Đạt; Đầy đủ: Thiếu ma trận toàn khóa; Tin cậy: Trung bình; Tương thích: Tốt. & \textbf{Đạt có điều kiện} \\
\textbf{TC 2.2} & Biên bản họp Hội đồng khoa học (\texttt{H1.02.02.01}); Nhận xét của 05 doanh nghiệp. & Pháp lý: Chuẩn; Đầy đủ: Khá; Tin cậy: Cao (có dấu doanh nghiệp); Tương thích: Tốt. & \textbf{Đạt yêu cầu} \\
\textbf{TC 4.2} & Hồ sơ bằng ĐH, ThS (\texttt{H1.04.02.01}); Chứng chỉ kỹ năng nghề Bậc 3/5 của 12 nhà giáo. & Pháp lý: Xuất sắc; Đầy đủ: Đạt 55\% tay nghề cao; Tin cậy: Rất cao; Tương thích: Tốt. & \textbf{Đạt yêu cầu} \\
\textbf{TC 5.3} & Bảng kiểm kê (\texttt{H1.05.01.01}); Nhật ký bảo dưỡng (\texttt{H1.05.03.01}); Biên bản tiếp địa. & Pháp lý: Đạt; Đầy đủ: Thiếu nhật ký máy dập phôi; Tin cậy: Trung bình; Tương thích: Đạt. & \textbf{Đạt có điều kiện} \\
\bottomrule
\end{tabularx}
\end{table}

Kết quả thẩm định độc lập cho thấy ba trên năm tiêu chuẩn đã đáp ứng xuất sắc các thuộc tính chất lượng của minh chứng, trong khi hai tiêu chuẩn (Tiêu chuẩn 1.3 và Tiêu chuẩn 5.3) cần phải được bổ sung và chuẩn hóa dữ liệu thực chứng trước khi bước vào đợt đánh giá ngoài chính thức.

\subsection{Báo cáo tự đánh giá hoàn chỉnh cho Tiêu chuẩn 1.3 về tích hợp năng lực số}

Nhằm minh họa quy chuẩn kỹ thuật viết báo cáo tự đánh giá theo đúng hướng dẫn của cơ quan kiểm định, Tiêu chuẩn 1.3 thuộc Tiêu chí 1 được biên soạn trọn vẹn với đầy đủ bốn cấu phần bắt buộc:

\subsubsection*{1. Mô tả hiện trạng}
Thực hiện chủ trương chuyển đổi số của Chính phủ theo Quyết định số 2222/QĐ-TTg \parencite{qd2222} và định hướng chỉ đạo của Bộ Lao động -- Thương binh và Xã hội, trong đợt rà soát chương trình đào tạo năm 2025, Khoa Điện -- Điện tử đã phối hợp cùng Hội đồng khoa học nhà trường tiến hành cập nhật, bổ sung các tiêu chuẩn năng lực số vào cấu trúc chuẩn đầu ra của chương trình đào tạo trình độ cao đẳng nghề Điện công nghiệp [\texttt{H1.01.01.01}].

Nội hàm năng lực số đã được cụ thể hóa thành các chỉ số hành vi đo lường được trong bản mô tả chuẩn đầu ra chính thức [\texttt{H1.01.02.01}]. Cụ thể, người học sau khi tốt nghiệp phải có khả năng sử dụng thành thạo các phần mềm chuyên ngành như AutoCAD Electrical và EPLAN Electric để thiết kế bản vẽ sơ đồ nguyên lý và bản vẽ lắp đặt hệ thống điện công nghiệp; có khả năng lập trình điều khiển, mô phỏng hoạt động của hệ thống tự động hóa sử dụng bộ điều khiển logic khả trình thông qua phần mềm TIA Portal; đồng thời có năng lực thu thập dữ liệu vận hành và giám sát trạng thái năng lượng của động cơ điện qua mạng công nghiệp Modbus TCP/IP [\texttt{H1.01.03.01}].

Các năng lực số này đã được tích hợp trực tiếp vào đề cương chi tiết của tám mô-đun chuyên ngành trọng điểm, chiếm 28.5\% tổng khối lượng mô-đun chuyên môn của toàn khóa học [\texttt{H1.01.03.02}]. Nhà trường đã đầu tư một phòng thực hành mô phỏng số chuyên dụng với ba mươi máy tính cấu hình cao kết nối mạng nội bộ và cài đặt đầy đủ các phần mềm kỹ thuật điện có bản quyền giáo dục để phục vụ cho sinh viên học tập và rèn luyện kỹ năng số trước khi thao tác trên thiết bị phần cứng thực tế [\texttt{H1.05.01.02}].

\subsubsection*{2. Điểm mạnh}
Nhà trường đã chủ động đi đầu trong việc xác lập khung năng lực số chi tiết cho nghề Điện công nghiệp, gắn liền năng lực số với các tác vụ kỹ thuật cụ thể mà các doanh nghiệp sản xuất trong khu công nghiệp đang đòi hỏi. Việc sinh viên được rèn luyện kỹ năng lập trình và thiết kế trên các phần mềm mô phỏng hiện đại giúp giảm thiểu nguy cơ tai nạn điện, tiết kiệm vật tư tiêu hao trong các giờ học ban đầu và nâng cao khả năng tiếp cận nhanh chóng với các dây chuyền sản xuất tự động hóa sau khi tốt nghiệp. Toàn bộ tài liệu hướng dẫn thực hành số hóa đã được đưa lên hệ thống học tập trực tuyến, tạo điều kiện thuận lợi cho sinh viên tự học và tự rèn luyện kỹ năng ngoài giờ lên lớp.

\subsubsection*{3. Tồn tại và nguyên nhân}
Mặc dù chuẩn đầu ra đã được ban hành và đưa vào đề cương mô-đun, việc triển khai đánh giá kết quả học tập đối với năng lực số vẫn chưa được tách biệt thành một cấu phần bài thi độc lập mà vẫn ghép chung vào điểm đánh giá thực hành truyền thống, khiến việc đo lường mức độ thuần thục kỹ năng số của từng sinh viên chưa thật sự sâu sắc. 

Nguyên nhân của hạn chế này là do một bộ phận nhà giáo lớn tuổi trong khoa còn lúng túng trong việc sử dụng các tính năng nâng cao của phần mềm mô phỏng công nghiệp; mặt khác, hệ thống phòng máy tính mô phỏng và xưởng thực hành phần cứng chưa được kết nối mạng truyền thông công nghiệp thời gian thực để truyền dữ liệu lập trình trực tiếp xuống các tủ điều khiển xưởng máy, gây ra sự phân tách nhất định giữa kỹ năng số trên máy tính và thao tác đấu nối phần cứng thực tế.

\subsubsection*{4. Kế hoạch cải tiến chất lượng}
Để khắc phục dứt điểm các tồn tại trên, nhà trường và Khoa Điện -- Điện tử xác lập kế hoạch hành động cải tiến chất lượng cụ thể:
\begin{itemize}
    \item Thứ nhất, tổ chức khóa đào tạo, bồi dưỡng nâng cao năng lực ứng dụng phần mềm thiết kế và điều khiển điện công nghiệp cho 100\% nhà giáo trong Khoa Điện -- Điện tử; thời gian hoàn thành trong quý IV năm 2026; đơn vị chủ trì là Khoa Điện -- Điện tử phối hợp cùng các chuyên gia kỹ thuật của hãng cung cấp thiết bị; kinh phí dự kiến trích từ quỹ đào tạo bồi dưỡng của trường là 120 triệu đồng.
    \item Thứ hai, hoàn thiện ma trận và bộ công cụ đánh giá xác thực năng lực số độc lập cho từng học phần chuyên môn, ban hành ngân hàng bài thi kỹ năng số trước tháng 12 năm 2026; người chịu trách nhiệm là Trưởng bộ môn Tự động hóa thuộc Khoa Điện -- Điện tử.
    \item Thứ ba, nâng cấp đường truyền mạng công nghiệp kết nối liên thông giữa phòng mô phỏng máy tính và ba xưởng thực hành điện để sinh viên có thể nạp chương trình điều khiển trực tiếp từ máy tính xuống mô hình vật lý tại xưởng; thời gian thực hiện trong quý I năm 2027; đơn vị chủ trì là Phòng Quản trị Thiết bị phối hợp Khoa Điện -- Điện tử.
    \item Chỉ số đầu ra cần đạt: 100\% sinh viên tốt nghiệp từ năm 2027 đạt chuẩn đầu ra năng lực số chuyên ngành được xác nhận thông qua bài kiểm tra đánh giá năng lực thực hành trên máy tính và dây chuyền mô phỏng.
\end{itemize}

\subsection{Điểm đánh giá dự kiến và đối chiếu các điều kiện công nhận đạt theo Điều 11 Thông tư số 14/2024/TT-BLĐTBXH}

Dựa trên kết quả tự rà soát toàn diện bốn mươi tiêu chuẩn kiểm định của chương trình đào tạo trình độ cao đẳng nghề Điện công nghiệp và căn cứ vào hồ sơ minh chứng đã được chuẩn bị, Hội đồng tự đánh giá nhà trường đã xây dựng bảng dự kiến điểm số chính thức trong Bảng~\ref{tab:bang2_4}.

\begin{table}[htbp]
\centering
\footnotesize
\setlength{\tabcolsep}{3.5pt}
\renewcommand{\arraystretch}{1.15}
\caption{\textbf{Bảng tổng hợp dự kiến điểm số tự đánh giá chương trình đào tạo nghề Điện công nghiệp}}
\label{tab:bang2_4}
\vspace{0.2cm}
\begin{tabularx}{\textwidth}{>{\centering\arraybackslash}p{1.1cm} X >{\centering\arraybackslash}p{0.9cm} >{\centering\arraybackslash}p{1.2cm} >{\centering\arraybackslash}p{1.0cm} >{\centering\arraybackslash}p{1.1cm} >{\centering\arraybackslash}p{1.1cm} >{\centering\arraybackslash}p{2.4cm}}
\toprule
\textbf{Mã TC} & \textbf{Tên tiêu chí kiểm định} & \textbf{Số TC} & \textbf{Điểm chuẩn} & \textbf{Số TC đạt} & \textbf{Điểm đạt} & \textbf{Tỷ lệ} & \textbf{Đánh giá theo Điều 11} \\
\midrule
TC-01 & Mục tiêu, chuẩn đầu ra & 3 & 7.5 & 3 & 7.5 & 100.0\% & Đạt ($\ge 50\%$) \\
TC-02 & Cấu trúc, nội dung, giáo trình & 7 & 17.5 & 6 & 15.0 & 85.7\% & Đạt ($\ge 75\%$) \\
TC-03 & Hoạt động đào tạo & 7 & 17.5 & 5 & 12.5 & 71.4\% & Đạt ($\ge 50\%$) \\
TC-04 & Đội ngũ nhà giáo, cán bộ quản lý & 6 & 15.0 & 5 & 12.5 & 83.3\% & Đạt ($\ge 75\%$) \\
TC-05 & Cơ sở vật chất, thiết bị, học liệu & 7 & 17.5 & 6 & 15.0 & 85.7\% & Đạt ($\ge 75\%$) \\
TC-06 & Người học và hoạt động hỗ trợ & 5 & 12.5 & 4 & 10.0 & 80.0\% & Đạt ($\ge 50\%$) \\
TC-07 & Giám sát, đánh giá chất lượng & 5 & 12.5 & 4 & 10.0 & 80.0\% & Đạt ($\ge 75\%$) \\
\midrule
\multicolumn{2}{l}{\textbf{Tổng số điểm đánh giá CTĐT}} & \textbf{40} & \textbf{100} & \textbf{33} & \textbf{82.5} & \textbf{82.5\%} & \textbf{ĐẠT CHUẨN} \\
\bottomrule
\end{tabularx}
\end{table}

Căn cứ theo quy định tại Khoản 1 Điều 11 Thông tư số 14/2024/TT-BLĐTBXH, chương trình đào tạo trình độ cao đẳng được đánh giá là đạt tiêu chuẩn kiểm định chất lượng khi đáp ứng đồng thời cả ba điều kiện pháp lý. Kết quả đối chiếu bảng điểm dự kiến với các quy định này được phân tích cụ thể:

Thứ nhất, đối chiếu với Điều kiện a về tổng điểm đánh giá, thông tư yêu cầu tổng số điểm đánh giá của chương trình đào tạo phải đạt từ 80 điểm trở lên. Điểm tự đánh giá dự kiến của chương trình đào tạo nghề Điện công nghiệp đạt 82.5 điểm trên 100 điểm, vượt ngưỡng sàn 2.5 điểm. Do đó, chương trình hoàn toàn đáp ứng được Điều kiện a.

Thứ hai, đối chiếu với Điều kiện b về ngưỡng chất lượng của các tiêu chí trọng tâm, thông tư quy định điểm đánh giá của các tiêu chí 2, 4, 5 và 7 phải đạt từ 75\% điểm chuẩn của từng tiêu chí trở lên. Kết quả tính toán cho thấy: Tiêu chí 2 đạt 15.0 trên 17.5 điểm (85.7\%); Tiêu chí 4 đạt 12.5 trên 15.0 điểm (83.3\%); Tiêu chí 5 đạt 15.0 trên 17.5 điểm (85.7\%); Tiêu chí 7 đạt 10.0 trên 12.5 điểm (80.0\%). Cả bốn tiêu chí cốt lõi này đều vượt trên mức sàn 75\% theo quy định. Do đó, chương trình hoàn toàn đáp ứng được Điều kiện b.

Thứ ba, đối chiếu với Điều kiện c về ngưỡng chất lượng đối với các tiêu chí thành phần còn lại, thông tư quy định điểm đánh giá của các tiêu chí 1, 3 và 6 phải đạt từ 50\% điểm chuẩn của từng tiêu chí trở lên. Kết quả tính toán ghi nhận: Tiêu chí 1 đạt 7.5 trên 7.5 điểm (100.0\%); Tiêu chí 3 đạt 12.5 trên 17.5 điểm (71.4\%); Tiêu chí 6 đạt 10.0 trên 12.5 điểm (80.0\%). Tất cả các tiêu chí này đều vượt trên mức tối thiểu 50\%. Do đó, chương trình hoàn toàn đáp ứng được Điều kiện c.

Mặc dù kết quả dự kiến đạt 82.5 điểm bảo đảm cho chương trình đào tạo nghề Điện công nghiệp đủ điều kiện pháp lý để được công nhận đạt tiêu chuẩn kiểm định, nhà trường cần đặc biệt lưu ý đến các rủi ro cận biên. Khoảng cách an toàn của tổng điểm so với ngưỡng sàn chỉ là 2.5 điểm (tương đương với điểm số của đúng một tiêu chuẩn). Nếu trong quá trình đánh giá ngoài thực tế, đoàn chuyên gia kiểm định độc lập bác bỏ tính xác thực của một tiêu chuẩn thuộc Tiêu chí 5 về an toàn xưởng hoặc Tiêu chí 7 về khảo sát việc làm, tổng điểm sẽ tụt xuống 80 điểm hoặc dưới 80 điểm, dẫn đến nguy cơ không đạt kiểm định. Do đó, việc xây dựng và vận hành một hệ thống đảm bảo chất lượng bên trong vững chắc để liên tục củng cố các tiêu chuẩn này là yêu cầu cấp thiết mang tính sống còn.
